\subsection{正则环与平坦扩张}

命题 8.--- 设 \(\rho: A \to B\) 为一个 Noether 环的同态，它使 B 成为忠实平坦 A-模。

\begin{enumerate}
\def\labelenumi{\alph{enumi})}
\item
  对每个有限型 \(\Lambda\) -模 M，有 \(\mathrm{dp}_{\mathrm{A}}(\mathrm{M}) = \mathrm{dp}_{\mathrm{B}}(\mathrm{B} \otimes_{\mathrm{A}} \mathrm{M})\) 。
\item
  若环 B 正则，则环 A 正则。
\end{enumerate}

B-模 \(B \otimes_{A} M\) 是有限型的 (I, §3, 第 6 节, 命题 12)。为使它为零，必须且只需 M 为零 (I, §3, 第 1 节, 定义 1)；为使它是投射的，必须且只需 M 是投射 A-模 (I, 第 6 节, 命题 12)。这就证明了 \(\mathrm{dp}_{\Lambda}(\mathrm{M}) \leqslant 0\) 情形的 a)。因此设 \(\mathrm{dp}_{\mathrm{A}}(\mathrm{M}) \geqslant 1\)（从而 \(\mathrm{dp}_{\mathrm{B}}(\mathrm{B} \otimes_{\Lambda} \mathrm{M}) \geqslant 1\)），并对 \(\mathrm{dp}_{\mathrm{A}}(\mathrm{M})\) 用归纳法证明 a)。选取一个 A-模的正合列

\[
0 \to \mathrm{N} \to \mathrm{L} \to \mathrm{M} \to 0
\]

其中 A-模 L 是有限型的自由模。我们有 \(\mathrm{dp}_{\Lambda}(\mathbf{N}) = \mathrm{dp}_{\mathbf{A}}(\mathbf{M}) - 1\)。由于 B 在 A 上平坦，序列

\[
0 \rightarrow \mathrm{B} \otimes_ {\mathrm{A}} \mathrm{N} \rightarrow \mathrm{B} \otimes_ {\mathrm{A}} \mathrm{L} \rightarrow \mathrm{B} \otimes_ {\mathrm{A}} \mathrm{M} \rightarrow 0
\]

是正合的，并且有 \(\mathrm{dp}_{\mathrm{B}}(\mathrm{B}\otimes_{\mathrm{A}}\mathrm{N}) = \mathrm{dp}_{\mathrm{B}}(\mathrm{B}\otimes_{\mathrm{A}}\mathrm{M}) - 1\)。把归纳假设应用于 N 即可得出结论。这就证明了 a)；断言 b) 由 a) 及第 2 节的命题 4 得出。

推论. 设 B 为正则整环，A 为 B 的 Noether 子环且 B 是有限生成 A-模。下列条件等价：

\begin{enumerate}
\def\labelenumi{(\roman{enumi})}
\item
  A 是正则的；
\item
  B 是投射 Λ-模；
\item
  B 是平坦 A-模；
\item
  B 是忠实平坦 A-模。
\item
  \(\Rightarrow\) (ii)：这由第 3 节的命题 5 的推论得出。
\item
  \(\Rightarrow\) (iii)：这由 I, § 3, 第 1 节, 命题 1 得出。
\item
  \(\Rightarrow\) (iv)：对 A 的每个素理想 \(\mathfrak{p}\)，有 \(\mathfrak{pB} \neq B\) (V, § 2, 第 1 节, 定理 1 的推论 1)。于是只需应用 I, § 3, 第 1 节, 命题 1。
\item
  \(\Rightarrow\) (i)：这由命题 8 b) 得出。
\end{enumerate}

对每个 Noether 局部环 \( A \)，用 \( \delta(A) \) 表示整数

\[
\delta (\mathrm{A}) = \left[ \mathfrak {m} _ {\mathrm{A}} / \mathfrak {m} _ {\mathrm{A}} ^ {2}: \kappa_ {\mathrm{A}} \right] - \dim (\mathrm{A}).
\]

回忆 (VIII, § 5, 第 1 节)，\(\delta(A)\) 总是非负的，而它的为零刻画了正则局部环。

设 \(\rho : A \to B\) 为 Noether 局部环的局部同态；由 \(\rho\) 导出一个 \(\kappa_{A}\) -线性的 \(\mathfrak{m}_{A}/\mathfrak{m}_{A}^{2} \longrightarrow \mathfrak{m}_{B}/\mathfrak{m}_{B}^{2}\)，从而导出一个 \(\kappa_{B}\) -线性同态

\[
d \rho : \kappa_ {B} \otimes_ {\kappa_ {A}} \left(\mathfrak {m} _ {A} / \mathfrak {m} _ {A} ^ {2}\right) \longrightarrow \mathfrak {m} _ {B} / \mathfrak {m} _ {B} ^ {2}.
\]

引理 1.--- 我们有

\[
\delta (\mathrm{B}) + \left[ \operatorname{Ker} (d \rho): \kappa_ {\mathrm{B}} \right] = \delta (\mathrm{A}) + \delta \left(\kappa_ {\mathrm{A}} \otimes_ {\mathrm{A}} ^ {\prime} \mathrm{B}\right) + (\dim (\mathrm{A}) - \dim (\mathrm{B}) + \dim \left(\kappa_ {\mathrm{A}} \otimes_ {\mathrm{A}} \mathrm{B}\right)).
\]

用 C 表示局部环 \(\kappa_{\mathrm{A}}\otimes_{\mathrm{A}}\mathrm{B}\) 。考察 B-模的正合列

\[
\mathrm{B} \otimes_ {\mathrm{A}} \mathfrak {m} _ {\mathrm{A}} \longrightarrow \mathfrak {m} _ {\mathrm{B}} \longrightarrow \mathfrak {m} _ {\mathrm{C}} \rightarrow 0;
\]

与 \(\kappa_{\mathrm{B}}\) 作张量积，得到一个 \(\kappa_{\mathrm{B}}\) -向量空间的正合列

\[
\kappa_ {\mathrm{B}} \otimes_ {\kappa_ {\Lambda}} (\mathfrak {m} _ {\mathrm{A}} / \mathfrak {m} _ {\mathrm{A}} ^ {2}) \stackrel {d \rho} {\longrightarrow} \mathfrak {m} _ {\mathrm{B}} / \mathfrak {m} _ {\mathrm{B}} ^ {2} \longrightarrow \mathfrak {m} _ {\mathrm{C}} / \mathfrak {m} _ {\mathrm{C}} ^ {2} \rightarrow 0,
\]

由此导出等式

\[
\left[ \mathfrak {m} _ {\mathrm{B}} / \mathfrak {m} _ {\mathrm{B}} ^ {2}: \kappa_ {\mathrm{B}} \right] + \left[ \operatorname{Ker} (d \boldsymbol {\rho}): \kappa_ {\mathrm{B}} \right] = \left[ \mathfrak {m} _ {\mathrm{A}} / \mathfrak {m} _ {\mathrm{A}} ^ {2}: \kappa_ {\mathrm{A}} \right] + \left[ \mathfrak {m} _ {\mathrm{C}} / \mathfrak {m} _ {\mathrm{C}} ^ {2}: \kappa_ {\mathrm{C}} \right],
\]

这就蕴含引理。

命题 9.- 设 \(\rho : A \to B\) 为 Noether 局部环的局部同态。下列条件等价：

\begin{enumerate}
\def\labelenumi{(\roman{enumi})}
\tightlist
\item
  环 B 正则，且 \(\kappa_{\mathrm{B}}\) -线性映射

\[
d \rho : \kappa_ {B} \otimes_ {\kappa_ {A}} \left(\mathfrak {m} _ {A} / \mathfrak {m} _ {A} ^ {2}\right) \longrightarrow \mathfrak {m} _ {B} / \mathfrak {m} _ {B} ^ {2}
\]

是单射的；

\item
  环 B 与 \(\kappa_{\mathrm{A}}\otimes_{\mathrm{A}}\mathrm{B}\) 都正则且 A-模 B 平坦；
\item
  环 A 与 \(\kappa_{\mathrm{A}}\otimes_{\mathrm{A}}\mathrm{B}\) 都正则且 A-模 B 平坦；
\item
  环 A 与 \(\kappa_{\Lambda} \otimes_{\mathrm{A}} \mathrm{B}\) 都正则，且有

\[
\dim (\mathrm{B}) = \dim (\mathrm{A}) + \dim (\kappa_ {\mathrm{A}} \otimes_ {\Lambda} \mathrm{B}).
\]
\end{enumerate}

我们有 \(\dim(B) \leqslant \dim(A) + \dim(\kappa_A \otimes_A B)\) (VIII, § 3, 第 4 节, 命题 7 的推论 1)；因此 (i) 与 (iv) 的等价性由引理 1 得出。在假设 (ii) 下，A-模 B 是忠实平坦的，因为有 \(\mathfrak{m}_A B \subset \mathfrak{m}_B \neq B\) (I, § 3, 第 1 节, 命题 1)，由命题 8 这蕴含 (iii)。蕴含关系 (iii)⇒(iv) 由 VIII 同上引文得出。

现在只需证明：当等价的条件 (i) 与 (iv) 满足时，A-模 B 是平坦的。设 x 为 A 的一个坐标系。由于 \(d\rho\) 是单射的，该序列在 \(\rho\) 下的像是 B 的一个坐标系的一部分。于是序列 x 对 \(\Lambda\) 与对 B 都是完全割的 (VIII, § 5, 第 3 节, 命题 2)，且生成 A 中的理想 \(\mathfrak{m}_{\Lambda}\)。根据 A, X, 第 159 页的注 3，A-模 \(\mathrm{Tor}_{1}^{\mathrm{A}}(\kappa_{\mathrm{A}},\mathrm{B})\) 同构于 \(\mathrm{H}_{1}(\mathbf{x},\mathrm{B})\)；从而它为零；由此得出 B 在 A 上平坦 (III, § 5, 第 2 节, 定理 1 及第 4 节, 命题 2)。

例.--- * 设 X, Y 为两个复解析簇，局部有限维，\(f\) 为从 X 到 Y 的一个态射，\(x\) 为 X 的一个点。考察局部同态 \(\rho: \mathcal{O}_{\mathrm{Y}, f(x)} \to \mathcal{O}_{\mathrm{X}, x}\)（与 \(f\) 相伴）。映射 \(d\rho\) 是切映射 \(\mathrm{T}_x(f): \mathrm{T}_x(\mathrm{X}) \to \mathrm{T}_{f(x)}(\mathrm{Y})\) 的转置。于是命题 9 的条件 (i) 至 (iv) 在此情形等价于 \(f\) 在 \(x\) 处是一个浸没 (VAR, R, 5.9.1)。*

推论.--- 设 \(\rho: \Lambda \to B\) 为一个 Noether 环的同态，它使 B 成为平坦 A-模。若 A 正则，且 \(\kappa(\rho^{-1}(\mathfrak{n})) \otimes_{\Lambda} B\) 对 B 的每个极大理想 \(\mathfrak{n}\) 都正则，则环 B 正则。

事实上，对 B 的每个极大理想 \(\mathfrak{n}\)，\(A_{\rho^{-1}(n)}\) -模 \(B_{n}\) 都是平坦的 (II, § 3, 第 4 节, 命题 15)，从而由命题 9，环 \(B_{n}\) 正则。

